\chapter{The Two-Torus: \texorpdfstring{$\tau$}{tau}, \texorpdfstring{$\rho$}{rho} and Mirror Symmetry}\label{ch:torus2}

A string on a circle has one modulus, the radius, and one duality, $R\to\ap/R$
(\cref{ch:tduality}). On a $d$-torus there are $d^2$ moduli and the duality group is
$\Odd{d}$ (\cref{ch:narain,ch:odd}). The first case in which this group is more than a
reflection is $d=2$, and it is the case a student should know by heart, for three reasons.
First, the four real moduli of $T^2$ combine into two complex numbers, $\tau$ and $\rho$,
each living in the upper half-plane, and the abstract group $\Odd{2}$ becomes two copies of
the modular group $\SLZ$ acting by fractional linear maps, which can be drawn. Second, the
circle duality applied to one side of the torus \emph{exchanges} $\tau$ and $\rho$: the shape
of one torus is the size of the other. This is the simplest example of mirror
symmetry\index{mirror symmetry}, and the only one in this book where everything can be
computed by hand. Third, the torus $T^2$ is a factor of the compactification space $\KT$ of
\cref{ch:k3t2}, and the hyperbolic geometry of its moduli space is the model for the distance
statements of \cref{ch:swampland}. The question of this chapter is therefore: \emph{what are
$\tau$ and $\rho$, how does each generator of $\Odd{2}$ act on them, and which of these
statements have been checked by the Lean kernel?}

\paragraph{Conventions.} As in \cref{ch:narain}, the torus coordinates have period $2\pi$,
$X^i\sim X^i+2\pi$, the metric $G_{ij}=\hat e_i\cdot\hat e_j/\ap$ and the antisymmetric field
$B_{ij}$ are dimensionless (lengths are measured in units of $\sqrt{\ap}$), the doubled
charge vector is $Z=(m^1,m^2,n_1,n_2)$ with windings first, and the generalized metric\index{generalized metric} is
\begin{equation}\label{eq:torus2:genmetric}
 \HH(G,B)=\begin{pmatrix}G-BG^{-1}B&BG^{-1}\\[2pt]-G^{-1}B&G^{-1}\end{pmatrix},
 \qquad \tfrac12(p_L^2+p_R^2)=\tfrac12 Z^{\mathsf T}\HH Z .
\end{equation}
A duality $g\in\Odd{2}$ acts on backgrounds by
\begin{equation}\label{eq:torus2:action}
 \HH\;\longmapsto\; g\,\HH\,g^{\mathsf T},
\end{equation}
which is the convention of our main source \source{GPR §2.4, eq.~(2.4.19), ll.~1288--1293}
and of the Lean theorem \lean{genMetric\_bshift}. The transposed convention
$\HH\mapsto g^{\mathsf T}\HH g$ of \cref{ch:narain} describes the same group; the difference
matters only when one asks which matrix does what, and we return to it in
\cref{sec:torus2:lean}.

\section{Four real numbers, two complex numbers}\label{sec:torus2:moduli}

\subsection{The shape of the torus: \texorpdfstring{$\tau$}{tau}}

A flat two-torus is the plane divided by the lattice generated by two vectors
$2\pi\hat e_1$, $2\pi\hat e_2$. Identify the plane with $\C$ and regard $\hat e_1,\hat e_2$ as
complex numbers. Rotating the plane and rescaling all lengths does not change the
\emph{shape} of the torus, that is, its structure as a Riemann surface; both operations
multiply $\hat e_1$ and $\hat e_2$ by the same non-zero complex number. The only invariant of
the pair under such a multiplication is the ratio
\begin{equation}\label{eq:torus2:taudef}
 \tau=\frac{\hat e_1}{\hat e_2}\in\C ,
\end{equation}
and we order the two vectors so that $\operatorname{Im}\tau>0$. This number is the
\emph{complex structure modulus}\index{complex structure modulus} of the torus. To express
it through the metric, multiply numerator and denominator by $\bar{\hat e}_2$:
$\tau=\hat e_1\bar{\hat e}_2/|\hat e_2|^2$. The real part of $\hat e_1\bar{\hat e}_2$ is the
scalar product $\hat e_1\cdot\hat e_2=\ap G_{12}$, and its imaginary part is, up to the sign
fixed by $\operatorname{Im}\tau>0$, the area of the parallelogram spanned by the two vectors,
$\ap\sqrt{\det G}$. Hence
\begin{equation}\label{eq:torus2:tau}
 \tau=\tau_1+\ii\tau_2=\frac{G_{12}}{G_{22}}+\ii\,\frac{\sqrt{\det G}}{G_{22}},
 \qquad \det G=G_{11}G_{22}-G_{12}^2>0
\end{equation}
\source{GPR §2.4.2, eq.~(2.4.55), ll.~1883--1892}. After a rotation and rescaling that bring
$\hat e_2$ to the number $1$, the torus is the parallelogram with corners $0,1,\tau,1+\tau$
(\cref{fig:torus2:lattice}). The angle $\theta$ between the two sides satisfies
$\cos\theta=\tau_1/|\tau|$, and the ratio of their lengths is $|\tau|$; a rectangular torus
has $\tau_1=0$, a square one $\tau=\ii$.

\begin{figure}[t]
\centering
\begin{tikzpicture}[scale=1.0,>=Stealth]
 \begin{scope}
 \foreach \a in {-1,0,1,2}{\foreach \b in {0,1,2}{
   \fill[black!50] ({1.6*\a+0.7*\b},{1.3*\b}) circle (1.2pt);}}
 \fill[blue!8] (0,0)--(1.6,0)--(2.3,1.3)--(0.7,1.3)--cycle;
 \draw[->,very thick,blue!60!black] (0,0)--(1.6,0) node[midway,below]{$\hat e_2\to1$};
 \draw[->,very thick,blue!60!black] (0,0)--(0.7,1.3) node[pos=0.75,left]{$\hat e_1\to\tau$};
 \draw[->,thick,red!70!black] (0,0)--(2.3,1.3);
 \node[red!70!black,right] at (2.3,1.35) {$\tau+1$};
 \draw[black!60] (0.45,0) arc (0:61.7:0.45); \node[black!70] at (0.62,0.3) {\small$\theta$};
 \node at (1.15,-0.95) {(a)};
 \end{scope}
 \begin{scope}[xshift=4.6cm]
 \draw[thick,fill=blue!8] (0,0) rectangle (3.2,1.2);
 \draw[<->] (0,-0.25)--(3.2,-0.25) node[midway,below]{\small$2\pi R_1$};
 \draw[<->] (3.45,0)--(3.45,1.2) node[midway,right]{\small$2\pi R_2$};
 \node[right] at (4.4,0.6) {\small$\tau=\ii\frac{R_1}{R_2},\ \ \rho=\ii\frac{R_1R_2}{\ap}$};
 \draw[->,thick] (1.6,1.45)--(1.6,2.15) node[midway,right]{\small$R_2\to\ap/R_2$};
 \draw[thick,fill=red!8] (0.72,2.4) rectangle (2.48,4.53);
 \node[right] at (4.4,3.45) {\small$\tau'=\rho,\ \ \rho'=\tau$};
 \node at (1.6,-0.95) {(b)};
 \end{scope}
\end{tikzpicture}
\caption{(a) The lattice of a two-torus, rotated and rescaled so that $\hat e_2=1$; then
$\hat e_1=\tau$. Replacing $\hat e_1$ by $\hat e_1+\hat e_2$ (red) gives the same lattice,
hence the same torus, with $\tau+1$ in place of $\tau$. (b) A rectangular torus with
$R_1/R_2=8/3$ and $R_1R_2=\frac32\ap$, and its image under T-duality on the second circle:
the ratio of the sides and the area in string units have been exchanged.}
\label{fig:torus2:lattice}
\end{figure}

\begin{pitfallbox}[Common pitfall: which ratio is $\tau$?]
Many texts define $\tau=\hat e_2/\hat e_1$, the inverse of \eqref{eq:torus2:taudef} up to
orientation. We follow GPR, where $G_{22}$ stands in the denominator of \eqref{eq:torus2:tau}.
With the other convention $\tau$ is replaced by $-1/\tau$ composed with a reflection, every
statement of this chapter survives, but the matrices that represent ``$\tau\to\tau+1$'' and
``$\tau\to\tau/(\tau+1)$'' exchange their roles. Always check a formula on the rectangular
torus before using it.
\end{pitfallbox}

\subsection{The size and the \texorpdfstring{$B$}{B}-field: \texorpdfstring{$\rho$}{rho}}

The metric has three independent components and $\tau$ uses two combinations of them. The
third is the overall scale, which we measure by the area. With coordinates of period $2\pi$
the area of the torus is $(2\pi)^2\ap\sqrt{\det G}$. In $d=2$ an antisymmetric matrix has one
component, $B_{12}=-B_{21}=b$. The \emph{complexified K\"ahler
modulus}\index{Kahler modulus@K\"ahler modulus} combines the two:
\begin{equation}\label{eq:torus2:rho}
 \rho=\rho_1+\ii\rho_2=B_{12}+\ii\sqrt{\det G}
\end{equation}
\source{GPR §2.4.2, eq.~(2.4.55), ll.~1883--1892}. The name comes from complex geometry:
on a complex manifold with a compatible metric, the K\"ahler form is the two-form built from
the metric and the complex structure; on $T^2$ it is the area form, and its integral is the
area. One half-plane of the $T^2$ moduli space ``represents the K\"ahler form and $B$-field
on the torus'' \source{Asp §5.1, ll.~2750--2762}. Why the area and $B$ belong together in one
complex number, rather than being two unrelated real parameters, is not obvious at this
point; the mass formula of \cref{sec:torus2:spectrum} will show that $\rho$ enters the
spectrum exactly as $\tau$ does.

The map $(G_{11},G_{12},G_{22},B_{12})\mapsto(\tau,\rho)$ is invertible. Solving
\eqref{eq:torus2:tau} and \eqref{eq:torus2:rho} for the metric gives
$G_{22}=\rho_2/\tau_2$, $G_{12}=\tau_1G_{22}$ and $G_{11}=(\det G+G_{12}^2)/G_{22}
=\rho_2(\tau_2^2+\tau_1^2)/\tau_2$, that is
\begin{equation}\label{eq:torus2:E}
 E=G+B=\frac{\rho_2}{\tau_2}\begin{pmatrix}|\tau|^2&\tau_1\\ \tau_1&1\end{pmatrix}
 +\rho_1\begin{pmatrix}0&1\\-1&0\end{pmatrix}
\end{equation}
\source{GPR §2.4.2, eq.~(2.4.56), ll.~1893--1917}. One checks
$\det G=(\rho_2/\tau_2)^2(|\tau|^2-\tau_1^2)=\rho_2^2$, as it must be. The moduli space of
backgrounds is therefore, before any duality is taken into account,
\begin{equation}\label{eq:torus2:HxH}
 \{(\tau,\rho)\}=\mathbb H\times\mathbb H,\qquad
 \mathbb H=\{z\in\C:\operatorname{Im}z>0\},
\end{equation}
two copies of the upper half-plane\index{upper half-plane} (the symbol $\mathbb H$ is not
to be confused with the generalized metric $\HH$). This is the $d=2$ case of the
coset $\OddR{2}/(O(2)\times O(2))$ of \cref{ch:narain}, and each factor is
$SL(2,\R)/U(1)$ \source{GPR §2.4.2, ll.~1876--1882; Asp §5.1, eq.~(100), ll.~2750--2762}.

\begin{example}[Rectangular torus]\label{ex:torus2:rect}
Two orthogonal circles of radii $R_1,R_2$ and $B=0$ have $G=\diag(R_1^2,R_2^2)/\ap$, so
\begin{equation}\label{eq:torus2:rect}
 \tau=\ii\,\frac{R_1}{R_2},\qquad \rho=\ii\,\frac{R_1R_2}{\ap}.
\end{equation}
The first is the ratio of the sides, the second the area in units of $(2\pi)^2\ap$.
T-duality on the second circle, $R_2\to\ap/R_2$, gives $\tau'=\ii R_1R_2/\ap=\rho$ and
$\rho'=\ii R_1/R_2=\tau$. For $R_1=2\sqrt{\ap}$, $R_2=\frac34\sqrt{\ap}$ one has
$\tau=\frac83\ii$, $\rho=\frac32\ii$ (\cref{fig:torus2:lattice}b): a long thin torus of
moderate area is dual to a moderately elongated torus of large area. This exchange is the
subject of \cref{sec:torus2:mirror}.
\end{example}

\section{The spectrum in terms of \texorpdfstring{$\tau$}{tau} and
\texorpdfstring{$\rho$}{rho}}\label{sec:torus2:spectrum}

The zero-mode spectrum of the string on $T^2$ is governed by the left and right momenta
$p_L,p_R$ of \cref{ch:narain}, whose squared lengths are
\begin{equation}\label{eq:torus2:pLpRdef}
 p_L^2=\tfrac12\,v_L^{\mathsf T}G^{-1}v_L,\quad v_L=n+(G-B)m;\qquad
 p_R^2=\tfrac12\,v_R^{\mathsf T}G^{-1}v_R,\quad v_R=n-(G+B)m .
\end{equation}
Here $n=(n_1,n_2)$ are the momentum numbers and $m=(m^1,m^2)$ the winding numbers; below we
write $m_i$ for $m^i$ to keep the formulas light. We now express \eqref{eq:torus2:pLpRdef}
through $\tau$ and $\rho$. The computation has three steps.

\paragraph{Step 1: the length of a covector.} Inverting the metric in \eqref{eq:torus2:E},
\[
 G^{-1}=\frac1{\rho_2\tau_2}\begin{pmatrix}1&-\tau_1\\-\tau_1&|\tau|^2\end{pmatrix},
 \qquad
 v^{\mathsf T}G^{-1}v=\frac{v_1^2-2\tau_1v_1v_2+|\tau|^2v_2^2}{\rho_2\tau_2}
 =\frac{|v_1-\tau v_2|^2}{\rho_2\tau_2}
\]
for any real $v=(v_1,v_2)$. The complex number $v_1-\tau v_2$ is, up to a phase and the
factor $(\rho_2\tau_2)^{-1/2}$, the covector $v$ written in an orthonormal frame.

\paragraph{Step 2: the map $v\mapsto v_1-\tau v_2$ applied to $Gm$ and $Bm$.} From
\eqref{eq:torus2:E},
\begin{align*}
 (Gm)_1-\tau(Gm)_2&=\frac{\rho_2}{\tau_2}\Big[|\tau|^2m_1+\tau_1m_2-\tau(\tau_1m_1+m_2)\Big]
 =\frac{\rho_2}{\tau_2}\Big[\tau(\bar\tau-\tau_1)\,m_1+(\tau_1-\tau)\,m_2\Big]\\
 &=\frac{\rho_2}{\tau_2}\,(-\ii\tau_2)\,(\tau m_1+m_2)=-\ii\rho_2\,(m_2+\tau m_1),
\end{align*}
where we used $|\tau|^2=\tau\bar\tau$ and $\bar\tau-\tau_1=\tau_1-\tau=-\ii\tau_2$. For the
$B$-field, $Bm=\rho_1(m_2,-m_1)$ and therefore $(Bm)_1-\tau(Bm)_2=\rho_1(m_2+\tau m_1)$. The
same complex combination $m_2+\tau m_1$ of the winding numbers appears in both.

\paragraph{Step 3: assemble.} For $v_L=n+Gm-Bm$ the coefficient of $(m_2+\tau m_1)$ is
$-\ii\rho_2-\rho_1=-\rho$; for $v_R=n-Gm-Bm$ it is $+\ii\rho_2-\rho_1=-\bar\rho$. Hence:

\begin{proposition}[Zero-mode spectrum on $T^2$]\label{thm:torus2:spectrum}
With $\tau,\rho$ defined by \eqref{eq:torus2:tau} and \eqref{eq:torus2:rho},
\begin{equation}\label{eq:torus2:pLpR}
 p_L^2=\frac{\big|(n_1-\tau n_2)-\rho\,(m_2+\tau m_1)\big|^2}{2\rho_2\tau_2},
 \qquad
 p_R^2=\frac{\big|(n_1-\tau n_2)-\bar\rho\,(m_2+\tau m_1)\big|^2}{2\rho_2\tau_2}.
\end{equation}
\end{proposition}

\noindent This is \source{GPR §2.4.2, eq.~(2.4.58), ll.~1953--1967}; we have also confirmed
by a symbolic check (sympy) that \eqref{eq:torus2:pLpR} equals \eqref{eq:torus2:pLpRdef}
identically in the eight variables, and that $p_L^2-p_R^2=2(n_1m_1+n_2m_2)$. The last
identity can be seen directly: with $a=n_1-\tau n_2$ and $c=m_2+\tau m_1$,
$|a-\rho c|^2-|a-\bar\rho c|^2=-2\operatorname{Re}\big[(\rho-\bar\rho)\,\bar a c\big]
=4\rho_2\operatorname{Im}(\bar ac)$ and
$\operatorname{Im}(\bar ac)=\tau_2(n_1m_1+n_2m_2)$, so the moduli cancel, as the
level-matching condition requires.

The mass formula of \cref{ch:narain} becomes
\begin{equation}\label{eq:torus2:mass}
 \frac\ap2M^2=\frac{p_L^2+p_R^2}{2}+N+\tilde N-2,\qquad N-\tilde N=n_1m_1+n_2m_2 .
\end{equation}

\begin{physicsbox}[Physical picture: what $\tau$ and $\rho$ do to the spectrum]
Set the windings to zero. Then $p_L^2=p_R^2=|n_1-\tau n_2|^2/(2\rho_2\tau_2)$: the
Kaluza--Klein masses are the squared lengths of the vectors of the \emph{dual} lattice,
whose shape is governed by $\tau$ and whose scale is $1/\rho_2$ --- a large torus has light
momentum modes. Now set the momenta to zero: $p_L^2=p_R^2=|\rho|^2\,|m_2+\tau m_1|^2/(2\rho_2\tau_2)$.
For $B=0$ this is $\rho_2|m_2+\tau m_1|^2/(2\tau_2)$, the squared length of the wound
string, which grows with the area. The $B$-field does not affect a state with momentum
only, but it shifts the effective momentum of a wound string: for $m\neq0$ the numerator
contains $n_1-\rho_1m_2$ and $n_2+\rho_1m_1$. A shift of $\rho_1$ by an integer can thus
be absorbed in a relabelling of the integers $n_i$.
\end{physicsbox}

The value of \eqref{eq:torus2:pLpR} is that symmetries can be read off. Any transformation
of $(\tau,\rho)$ that can be compensated by an invertible integer relabelling of
$(n_1,n_2,m_1,m_2)$, leaving the set of pairs $(p_L^2,p_R^2)$ unchanged, maps a background to
a physically equivalent one. For example, exchanging $\tau\leftrightarrow\rho$ in the
numerator of $p_L^2$ gives $n_1-\rho n_2-\tau m_2-\tau\rho m_1$, which is the original
expression $n_1-\tau n_2-\rho m_2-\rho\tau m_1$ with $n_2$ and $m_2$ exchanged: a momentum
number has been traded for a winding number in the second direction, which is what
T-duality on the second circle does. We now go through the generators systematically.

\section{The duality group of the two-torus}\label{sec:torus2:group}

We use the three families of generators of $\Odd{d}$ from \cref{ch:odd}
\source{GPR §2.4, eqs.~(2.4.25), (2.4.26), (2.4.29), ll.~1355--1414}, in block form with
respect to $Z=(m,n)$:
\begin{equation}\label{eq:torus2:gens}
 g_\Theta=\begin{pmatrix}\mathbf 1&\Theta\\0&\mathbf 1\end{pmatrix},\qquad
 g_A=\begin{pmatrix}A&0\\0&(A^{\mathsf T})^{-1}\end{pmatrix},\qquad
 g_{D_i}=\begin{pmatrix}\mathbf 1-e_i&e_i\\e_i&\mathbf 1-e_i\end{pmatrix},
\end{equation}
with $\Theta$ an antisymmetric integer matrix, $A\in GL(2,\Z)$, and $e_i$ the matrix with a
single $1$ at position $(i,i)$. In Lean these are \lean{thetaShift}, \lean{basisChange} and
\lean{factorized}. Inserting them in \eqref{eq:torus2:action} and comparing blocks with
\eqref{eq:torus2:genmetric} gives their action on $E=G+B$:
\begin{equation}\label{eq:torus2:Eaction}
 g_\Theta:\ B\to B+\Theta,\qquad g_A:\ G\to AGA^{\mathsf T},\ \ B\to ABA^{\mathsf T}.
\end{equation}
For the first, this is the content of \lean{genMetric\_bshift} \TA. For the second,
$g_A\HH g_A^{\mathsf T}$ has lower right block $A^{-\mathsf T}G^{-1}A^{-1}=(AGA^{\mathsf T})^{-1}$
and upper right block $ABG^{-1}A^{-1}=(ABA^{\mathsf T})(AGA^{\mathsf T})^{-1}$.

\subsection{The modular group of the shape: \texorpdfstring{$SL(2,\Z)_\tau$}{SL(2,Z) tau}}

Let $A=\big(\begin{smallmatrix}a&b\\c&d\end{smallmatrix}\big)$ have integer entries and
$\det A=1$. The transformation $G\to AGA^{\mathsf T}$ is what the metric
$G_{ij}=\hat e_i\cdot\hat e_j/\ap$ undergoes when the lattice basis is changed to
\begin{equation}\label{eq:torus2:basis}
 \hat e_1'=a\,\hat e_1+b\,\hat e_2,\qquad \hat e_2'=c\,\hat e_1+d\,\hat e_2 .
\end{equation}
Since $\det A=1$ and the inverse matrix is again integral, the new vectors generate the
same lattice: the torus has not changed, only its description. Dividing the two equations
and using \eqref{eq:torus2:taudef},
\begin{equation}\label{eq:torus2:mobius}
 \tau'=\frac{\hat e_1'}{\hat e_2'}=\frac{a\tau+b}{c\tau+d} .
\end{equation}
What happens to $\rho$? The area of the cell is unchanged because $\det(AGA^{\mathsf T})=
(\det A)^2\det G$. For the $B$-field we need a small lemma, which is special to $2\times2$
matrices: with $J=\big(\begin{smallmatrix}0&1\\-1&0\end{smallmatrix}\big)$,
\begin{equation}\label{eq:torus2:AJAT}
 A\,J\,A^{\mathsf T}=(\det A)\,J\qquad\text{for every }2\times2\text{ matrix }A .
\end{equation}
Indeed the $(1,2)$ entry of the left side is $ad-bc$, the diagonal entries are $ab-ba=0$ and
$cd-dc=0$, and the result is antisymmetric. This is the Lean theorem
\lean{mul\_jMat\_mul\_transpose} \TA. Since $B=\rho_1J$, a basis change with $\det A=1$ leaves
$B$, hence $\rho$, untouched:
\begin{equation}\label{eq:torus2:SLtau}
 SL(2,\Z)_\tau:\qquad \tau\to\frac{a\tau+b}{c\tau+d},\qquad\rho\to\rho .
\end{equation}
The group is generated by
$T=\big(\begin{smallmatrix}1&1\\0&1\end{smallmatrix}\big)$, $\tau\to\tau+1$, and
$S=\big(\begin{smallmatrix}0&1\\-1&0\end{smallmatrix}\big)$, $\tau\to-1/\tau$ (the fact that
$S$ and $T$ generate $\SLZ$ is standard; not checked against a source in this book's
library). Geometrically $T$ is the change of basis drawn in \cref{fig:torus2:lattice}a and
$S$ exchanges the two sides of the cell, with a sign to preserve the orientation. These
symmetries ``reflect the fact that the target space is a 2-d torus, and thus unmodified by
modular transformations of its complex structure'' \source{GPR §2.4.2, ll.~1948--1949}: they
are present for a point particle on $T^2$ as well and are not stringy.

A basis change with $\det A=-1$ reverses the orientation. For the reflection
$R_1=\diag(-1,1)$, that is $X^1\to-X^1$, one finds $G_{12}\to-G_{12}$ and, by
\eqref{eq:torus2:AJAT}, $B\to-B$:
\begin{equation}\label{eq:torus2:R}
 R:\qquad(\tau,\rho)\to(-\bar\tau,-\bar\rho)
\end{equation}
\source{GPR §2.4.2, eq.~(2.4.60), ll.~1980--1986}.

\subsection{The modular group of the size: \texorpdfstring{$SL(2,\Z)_\rho$}{SL(2,Z) rho}}

In $d=2$ the antisymmetric integer matrices are $\Theta=tJ$, $t\in\Z$, and by
\eqref{eq:torus2:Eaction}
\begin{equation}\label{eq:torus2:Tprime}
 T'^{\,t}=g_{tJ}:\qquad \rho\to\rho+t,\qquad \tau\to\tau .
\end{equation}
The reason this is a symmetry\index{B-field@$B$-field!integer shift} was given in \cref{ch:odd}: the $B$-term of the worldsheet
action is a total derivative, it only measures how many times the worldsheet wraps the
torus, and an integer shift of $B_{12}$ changes the action by a multiple of $2\pi$
\source{GPR §2.4, ll.~1368--1371}. It ``expresses the periodicity in the antisymmetric
torsion $B_{12}$'' \source{GPR §2.4.2, l.~1950}. The $B$-field is thus an angle, like the
$\theta$-parameter of a gauge theory, and $\rho_1$ and $\tau_1$ are both periodic variables.

The analogy between \eqref{eq:torus2:Tprime} and $\tau\to\tau+1$ suggests that there is also
a symmetry
\begin{equation}\label{eq:torus2:Sprime}
 S':\qquad\rho\to-\frac1\rho,\qquad\tau\to\tau ,
\end{equation}
and \eqref{eq:torus2:pLpR} confirms it. Write $a=n_1-\tau n_2$ and $c=m_2+\tau m_1$, so that
$p_L^2=|a-\rho c|^2/(2\rho_2\tau_2)$. Under $\rho\to-1/\rho$ the imaginary part becomes
$\rho_2/|\rho|^2$, and a state with new charges $a',c'$ has
\[
 p_L'^{\,2}=\frac{|\rho|^2\,\big|a'+c'/\rho\big|^2}{2\rho_2\tau_2}
 =\frac{|\rho a'+c'|^2}{2\rho_2\tau_2},
\]
which equals the old $p_L^2$ for $a'=c$, $c'=-a$; the same substitution works for $p_R^2$
with $\bar\rho$. In terms of the integers, $a'=c$ means $(n_1',n_2')=(m_2,-m_1)$ and $c'=-a$
means $(m_1',m_2')=(n_2,-n_1)$: momenta and windings are exchanged, with a rotation by
$90^\circ$ in the lattice. For $B=0$ it sends the area $\rho_2$ to
$1/\rho_2$ at fixed shape. This is ``a stringy symmetry that for $B=0$ takes the volume in
target space into its inverse'' \source{GPR §2.4.2, ll.~1950--1952}, the two-dimensional
version of $R\to\ap/R$. Together,
\begin{equation}\label{eq:torus2:SLrho}
 SL(2,\Z)_\rho:\qquad \rho\to\frac{a'\rho+b'}{c'\rho+d'},\qquad\tau\to\tau .
\end{equation}

The two groups act on different variables, so they commute as transformations of
$\mathbb H\times\mathbb H$. At the level of the $4\times4$ integer matrices this is a
statement that needs proof. For the translations it follows from \eqref{eq:torus2:AJAT}:
\begin{equation}\label{eq:torus2:commute}
 g_A\,g_{tJ}=\begin{pmatrix}A&tAJ\\0&A^{-\mathsf T}\end{pmatrix},\qquad
 g_{tJ}\,g_A=\begin{pmatrix}A&tJA^{-\mathsf T}\\0&A^{-\mathsf T}\end{pmatrix},
\end{equation}
and $AJ=JA^{-\mathsf T}$ is equivalent to $AJA^{\mathsf T}=J$, that is, to $\det A=1$. This
is the Lean theorem \lean{basisChange\_comm\_thetaShift} \TA{} (\cref{sec:torus2:lean}).

\subsection{The exchanges: factorized duality and the full group}

The factorized dualities\index{factorized duality}\index{T-duality} $D_1,D_2$ of \eqref{eq:torus2:gens} are T-duality on the first and
on the second circle. Their action on the moduli, obtained by inserting
\eqref{eq:torus2:E} in \eqref{eq:torus2:action} and reading off the new $G$ and $B$ (a symbolic
check; the case $D_2$ is \source{GPR §2.4.2, eq.~(2.4.59), ll.~1969--1978}), is
\begin{equation}\label{eq:torus2:D}
 D_2:\ (\tau,\rho)\to(\rho,\tau),\qquad
 D_1:\ (\tau,\rho)\to\Big(-\frac1\rho,-\frac1\tau\Big),\qquad
 D_1D_2=\eta:\ (\tau,\rho)\to\Big(-\frac1\tau,-\frac1\rho\Big).
\end{equation}
The last entry is the inversion $E\to E^{-1}$ of \cref{ch:odd}; on $T^2$ it is the product
$SS'$ of the two modular inversions. The matrix identity $D_1D_2=\eta$ is
\lean{factorized\_two} \TA. The asymmetry between $D_1$ and $D_2$ is a consequence of the
convention \eqref{eq:torus2:taudef}, in which $\hat e_2$ is the reference side. As integer
matrices $D_1=-\,g_S\,S'D_2$ (symbolic check), so both dualities exchange the complex structure
with the K\"ahler modulus, and they differ by an element of
$SL(2,\Z)_\tau\times SL(2,\Z)_\rho$.

Finally, worldsheet parity\index{worldsheet parity} $\sigma\to-\sigma$ reverses the sign of $B$ and nothing else:
\begin{equation}\label{eq:torus2:W}
 W:\qquad(\tau,\rho)\to(\tau,-\bar\rho).
\end{equation}
It exchanges $p_L$ and $p_R$, changes the sign of $p_L^2-p_R^2$, and is therefore \emph{not}
an element of $\Odd{2}$, although it is a symmetry of the bosonic string
\source{GPR §2.4.2, eq.~(2.4.61), ll.~1986--1996}.

\begin{theorem}[Duality group of $T^2$ \TL]\label{thm:torus2:group}\index{duality group!of the two-torus}
The elements $S,T$ (acting on $\tau$), $S',T'$ (acting on $\rho$), the exchange $D_2$ and the
reflection $R$ generate $\Odd{2}$. The subgroup generated by $S,T,S',T'$ is normal, and the
quotient is $\Z_2\times\Z_2$, generated by the classes of $D_2$ and $R$:
\begin{equation}\label{eq:torus2:structure}
 \Odd{2}\;\cong\;\big[SL(2,\Z)_\tau\times SL(2,\Z)_\rho\big]\rtimes\big(\Z_2\times\Z_2\big),
\end{equation}
where the product in brackets is understood with $(-\mathbf 1,-\mathbf 1)$ identified with
the identity (\cref{rem:torus2:kernel}). Including $W$, a minimal set of generators of the full
symmetry group is $S,T,D_2,W$, with
$S'=D_2SD_2$, $T'=D_2TD_2$ and $R=D_2WD_2W$.
\end{theorem}

\noindent\source{GPR §2.4.2, ll.~1876--1882 and 1997--2000; generation of $\Odd{d}$ by
\eqref{eq:torus2:gens}: GPR §2.4, ll.~1422--1423}. GPR write the right-hand side as
``$SL(2,\Z)\times SL(2,\Z)\otimes_S[\Z_2\times\Z_2]$'' and stress that it is \emph{not} the
naive factorization into two modular groups that the local form
$\OddR{2}\sim SL(2,\R)\times SL(2,\R)$ might suggest. The relations $T'=D_2TD_2$ and
$S'=D_2SD_2$ are immediate from \eqref{eq:torus2:D}, since $D_2$ simply renames the variable
on which $T$ or $S$ acts; as integer matrices we have checked $D_2\,g_T\,D_2=g_J$ symbolically,
and that $D_2$ and $R$ commute, so that they generate a $\Z_2\times\Z_2$.

\begin{remark}[The element $-\mathbf 1$]\label{rem:torus2:kernel}
The matrix $-\mathbf 1_2\in SL(2,\Z)_\tau$ acts trivially on $\tau$ but not on the charges:
$g_{-\mathbf 1}=-\mathbf 1_4$ sends $Z\to-Z$. The same is true of $S'^2$: from $S'=D_2SD_2$
and $g_S^2=-\mathbf 1_4$ one gets $S'^2=-\mathbf 1_4$ (symbolic check). Thus the two centres
are mapped to the \emph{same} element of $\Odd{2}$, and the subgroup generated by
$S,T,S',T'$ is $\big(SL(2,\Z)\times SL(2,\Z)\big)/\Z_2$. On the moduli space only
$PSL(2,\Z)=\SLZ/\{\pm\mathbf 1\}$ acts effectively in each factor. The state
$Z$ and the state $-Z$ are a particle and its antiparticle; $-\mathbf 1_4$ is charge
conjugation.
\end{remark}

\begin{pitfallbox}[Common pitfall: ``$O(2,2)$ is $SL(2)\times SL(2)$''] 
Over the real numbers and for the connected component this is true, and it is why the
moduli space is a product of two half-planes. Over the integers, the two factors are
exchanged by an element of the group itself. Consequently the moduli space of the string on
$T^2$ is \emph{not} the product of two modular curves, but its quotient by the exchange
$\tau\leftrightarrow\rho$ and by $(\tau,\rho)\to(-\bar\tau,-\bar\rho)$: a torus of shape
$\tau$ and size $\rho$ is the same string background as one of shape $\rho$ and size $\tau$.
\end{pitfallbox}

\section{The upper half-plane}\label{sec:torus2:halfplane}

Both moduli live in $\mathbb H$ and both are acted on by $\SLZ$ through fractional linear
maps. This section collects the geometry of that action. Nothing in it is specific to
strings; the same half-plane appears as the moduli space of the worldsheet torus in
\cref{ch:narain} and as that of the type IIB axio-dilaton in \cref{ch:stabilization}. Except where a source is
given, the statements of this section are standard and derived on the page; they have not
been checked against a source in this book's library.

\subsection{Fractional linear maps}

For a real matrix $\gamma=\big(\begin{smallmatrix}a&b\\c&d\end{smallmatrix}\big)$ with
$ad-bc=1$ and $z\in\mathbb H$ put $\gamma z=(az+b)/(cz+d)$\index{Mobius transformation@M\"obius transformation}. Multiplying numerator and
denominator by $c\bar z+d$,
\begin{equation}\label{eq:torus2:Im}
 \gamma z=\frac{ac|z|^2+bd+adz+bc\bar z}{|cz+d|^2}
 \quad\Longrightarrow\quad
 \operatorname{Im}(\gamma z)=\frac{(ad-bc)\operatorname{Im}z}{|cz+d|^2}
 =\frac{\operatorname{Im}z}{|cz+d|^2}>0 ,
\end{equation}
so $\gamma$ maps $\mathbb H$ to itself; the denominator cannot vanish because $z$ is not
real. Two facts make this an action of the group: the identity matrix acts trivially, and
\begin{equation}\label{eq:torus2:compose}
 \gamma(\gamma'z)=(\gamma\gamma')z .
\end{equation}
The second is a short computation with fractions, and is kernel-checked for complex
matrices as \lean{mobius\_compose} \TA{} (\cref{sec:torus2:lean}). Since $\gamma$ and $-\gamma$
give the same map, the group acting effectively is $PSL(2,\R)$. The action is transitive:
$z=x+\ii y$ is the image of $\ii$ under
$\big(\begin{smallmatrix}\sqrt y&x/\sqrt y\\0&1/\sqrt y\end{smallmatrix}\big)$. The matrices
fixing $\ii$ satisfy $a\ii+b=\ii(c\ii+d)$, that is $a=d$, $b=-c$, $a^2+b^2=1$: they form the
rotation group $SO(2)=U(1)$. This proves $\mathbb H=SL(2,\R)/U(1)$, the statement used in
\eqref{eq:torus2:HxH}.

\subsection{The fundamental domain}

Two points of $\mathbb H$ related by an element of the modular group\index{modular group} $\SLZ$ describe the same torus. A
\emph{fundamental domain}\index{fundamental domain} is a region containing exactly one
representative of each orbit, up to boundary identifications. The standard choice is
\begin{equation}\label{eq:torus2:F}
 \mathcal F=\big\{z\in\mathbb H:\ |\operatorname{Re}z|\le\tfrac12,\ |z|\ge1\big\},
\end{equation}
drawn in \cref{fig:torus2:fd}. That every $z$ can be brought into $\mathcal F$ is shown by
an algorithm which is also the practical way to reduce a given modulus:
\begin{enumerate}[nosep]
\item apply a power of $T$ to achieve $|\operatorname{Re}z|\le\frac12$;
\item if $|z|<1$, apply $S$: by \eqref{eq:torus2:Im} this \emph{increases} the imaginary
part, $\operatorname{Im}(-1/z)=\operatorname{Im}z/|z|^2$; return to step 1.
\end{enumerate}
The algorithm terminates. Indeed, all values taken by the imaginary part are of the form
$\operatorname{Im}z/|cz+d|^2$ with integers $(c,d)$, and only finitely many pairs of integers
satisfy $|cz+d|<1$, because the numbers $cz+d$ form a lattice in $\C$. A strictly
increasing sequence in a set with finitely many values above its starting point must stop.
In Lean's mathematical library this is the theorem \lean{ModularGroup.exists\_smul\_mem\_fd}
of Mathlib (file \texttt{Mathlib/NumberTheory/Modular.lean}); it is part of the library this
project builds on, not a result of this project. That no two interior points of
$\mathcal F$ are equivalent is the companion statement
\lean{ModularGroup.eq\_smul\_self\_of\_mem\_fdo\_mem\_fdo}.

\begin{figure}[t]
\centering
\begin{tikzpicture}[scale=2.3,>=Stealth]
 \fill[blue!10] (-0.5,2.55)--(-0.5,0.866) arc (120:60:1)--(0.5,2.55)--cycle;
 \draw[->] (-1.75,0)--(1.85,0) node[right]{\small$\operatorname{Re}z$};
 \draw[->] (0,0)--(0,2.7) node[above]{\small$\operatorname{Im}z$};
 \draw[black!45] (-1.5,0)--(-1.5,2.55) (1.5,0)--(1.5,2.55);
 \begin{scope}
 \clip (-1.7,0) rectangle (1.7,2.55);
 \foreach \c in {-2,-1,0,1,2}{\draw[black!45] (\c+1,0) arc (0:180:1);}
 \foreach \c in {-5,-4,-2,-1,1,2,4,5}{\draw[black!30] (\c/3+1/3,0) arc (0:180:1/3);}
 \draw[black!45] (-0.5,0)--(-0.5,2.55) (0.5,0)--(0.5,2.55);
 \end{scope}
 \draw[very thick,blue!60!black] (-0.5,2.55)--(-0.5,0.866) arc (120:60:1)--(0.5,2.55);
 \fill (0,1) circle (0.9pt) node[above left]{\small$\ii$};
 \fill (0.5,0.866) circle (0.9pt) node[right]{\small$\ee^{\ii\pi/3}$};
 \fill (-0.5,0.866) circle (0.9pt) node[left]{\small$\ee^{2\pi\ii/3}$};
 \node at (0.22,1.95) {$\mathcal F$};
 \node[black!60] at (1.0,1.95) {\small$T\mathcal F$};
 \node[black!60] at (-1.0,1.95) {\small$T^{-1}\mathcal F$};
 \node[black!60] at (-0.2,0.8) {\small$S\mathcal F$};
 \draw[red!70!black,thick] (0,1)--(0,2.55);
 \draw[red!70!black,thick,dashed] (1.3,0) arc (0:180:0.9);
 \foreach \x/\l in {-1.5/{-\frac32},-1/{-1},-0.5/{-\frac12},0.5/{\frac12},1/{1},1.5/{\frac32}}{
   \draw (\x,0.03)--(\x,-0.03) node[below]{\scriptsize$\l$};}
\end{tikzpicture}
\caption{The fundamental domain $\mathcal F$ of $\SLZ$ in the upper half-plane (shaded) and
some of its images. The vertical sides are identified by $T$, the two halves of the arc by
$S$. The marked points have non-trivial stabilizers. Red: two geodesics of the Poincar\'e
metric, a vertical line (solid) and a semicircle centred on the real axis (dashed).}
\label{fig:torus2:fd}
\end{figure}

Three points of $\mathcal F$ are special because they are fixed by a non-trivial element of
$PSL(2,\Z)$: $z=\ii$ is fixed by $S$ (order 2), $z=\ee^{\ii\pi/3}$ by
$TS:z\to1-1/z$ (order 3), and its $T$-image $\ee^{2\pi\ii/3}$ by $ST$. At these points the
quotient $\SLZ\backslash\mathbb H$ is not smooth but has a conical (``orbifold'') singularity.
For $\tau$, these are the square and the hexagonal lattice, whose rotation symmetry is
$\Z_4$ and $\Z_6$ instead of the generic $\Z_2$. In the string moduli space the two points
\begin{equation}\label{eq:torus2:special}
 (\tau,\rho)=(\ii,\ii)\qquad\text{and}\qquad
 (\tau,\rho)=\big(\ee^{\ii\pi/3},\ee^{\ii\pi/3}\big)
\end{equation}
are the backgrounds with gauge symmetry $(SU(2)\times SU(2))_L\times(SU(2)\times SU(2))_R$
and $SU(3)_L\times SU(3)_R$ respectively \source{GPR §2.4.2, eqs.~(2.4.63)--(2.4.65),
ll.~2019--2052}; the second was studied in \cref{ch:narain}. Orbifold points\index{enhanced gauge symmetry} of the moduli
space and enhanced gauge symmetry go together: the extra massless vectors exist
because a finite duality group maps the background to itself.

\subsection{The Poincar\'e metric}

Is there a natural notion of distance between two tori? A metric on $\mathbb H$ that does
not distinguish between equivalent descriptions must be invariant under $\SLZ$; we show that
there is a metric invariant under all of $SL(2,\R)$ and that it is unique up to a constant.
Differentiating $\gamma z$ gives
$\dd(\gamma z)=\big[a(cz+d)-c(az+b)\big]\dd z/(cz+d)^2=\dd z/(cz+d)^2$. Combined with
\eqref{eq:torus2:Im},
\begin{equation}\label{eq:torus2:poincare}
 \frac{|\dd(\gamma z)|^2}{(\operatorname{Im}\gamma z)^2}
 =\frac{|\dd z|^2}{|cz+d|^4}\cdot\frac{|cz+d|^4}{(\operatorname{Im}z)^2}
 \qquad\Longrightarrow\qquad
 \dd s^2=\frac{\dd x^2+\dd y^2}{y^2}\quad\text{is invariant.}
\end{equation}
This is the \emph{Poincar\'e metric}\index{Poincar\'e metric}. Uniqueness: an invariant
metric is determined by its value at $\ii$ because the action is transitive, and at $\ii$ it
must be invariant under the rotations that fix $\ii$, hence proportional to
$\dd x^2+\dd y^2$. Its Gaussian curvature follows from the formula
$K=-\ee^{-2u}(\partial_x^2+\partial_y^2)u$ for a metric $\ee^{2u}(\dd x^2+\dd y^2)$: with
$u=-\ln y$ one has $(\partial_x^2+\partial_y^2)u=1/y^2$ and
\begin{equation}\label{eq:torus2:K}
 K=-y^2\cdot\frac1{y^2}=-1 .
\end{equation}
The half-plane is the hyperbolic plane\index{hyperbolic plane}, of constant negative
curvature. The fundamental domain, although it extends to $y=\infty$, has finite area:
\begin{equation}\label{eq:torus2:area}
 \operatorname{Area}(\mathcal F)=\int_{-1/2}^{1/2}\!\dd x\int_{\sqrt{1-x^2}}^{\infty}\frac{\dd y}{y^2}
 =\int_{-1/2}^{1/2}\frac{\dd x}{\sqrt{1-x^2}}=2\arcsin\tfrac12=\frac\pi3 .
\end{equation}

On the string moduli space $\mathbb H\times\mathbb H$ the corresponding metric is
\begin{equation}\label{eq:torus2:modmetric}
 \dd s^2=\frac{|\dd\tau|^2}{\tau_2^2}+\frac{|\dd\rho|^2}{\rho_2^2}
 =-\frac14\Tr\big(\dd\HH\,\dd\HH^{-1}\big),
\end{equation}
where the second equality is a symbolic check (sympy) using \eqref{eq:torus2:E}. The right
side is manifestly invariant under $\HH\to g\HH g^{\mathsf T}$ for constant $g$, and it
treats $\tau$ and $\rho$ symmetrically, as mirror symmetry demands. The expression is not
arbitrary: the action of double field theory contains the term
$\frac18\HH^{MN}\partial_M\HH^{KL}\partial_N\HH_{KL}$
\source{HHZ §1, eq.~(1.15), ll.~334--346}, which for moduli depending only on the
non-compact coordinates $x^\mu$ is proportional to the scalar kinetic term
$g_{ij}\,\partial_\mu\phi^i\partial^\mu\phi^j$, with $\phi^i=(\tau_1,\tau_2,\rho_1,\rho_2)$
and $g_{ij}$ the metric \eqref{eq:torus2:modmetric} (\cref{ch:dftaction}). The reduction itself is not carried out
in this book; the normalization of \eqref{eq:torus2:modmetric} is a convention. For the
circle, $\HH=\diag(R^2/\ap,\ap/R^2)$ gives $\dd s^2=2(\dd R/R)^2$, which is invariant under
$R\to\ap/R$.

\subsection{Geodesics and distances}

The geodesics\index{geodesic} of \eqref{eq:torus2:poincare} are the paths $(x(s),y(s))$
extremizing $\int L\,\dd s$ with $L=\frac12(\dot x^2+\dot y^2)/y^2$. Rather than writing
the Euler--Lagrange equations we use two conserved quantities. $L$ does not depend on $x$,
so the momentum conjugate to $x$ is conserved; $L$ does not depend on $s$, so the
``energy'', here equal to $L$ itself, is conserved:
\begin{equation}\label{eq:torus2:conserved}
 p=\frac{\dot x}{y^2}=\text{const},\qquad
 \frac{\dot x^2+\dot y^2}{y^2}=1 ,
\end{equation}
where we fixed the constant by choosing $s$ to be the arc length. The energy is a sum of two
non-negative terms; that $(\dot y/y)^2\ge0$ is the content of the (elementary) Lean lemma
\lean{poincare\_geodesic\_kinetic\_energy\_nonneg}.

\emph{Case $p=0$.} Then $x$ is constant and $\dot y/y=\pm1$: the geodesic is the vertical
line $y=y_0\ee^{\pm s}$. The distance between $\ii y_0$ and $\ii y_1$ is $|\ln(y_1/y_0)|$.

\emph{Case $p\ne0$.} Substituting $\dot x=py^2$ in the energy gives
$\dot y^2=y^2(1-p^2y^2)$, so
$\dd x/\dd y=\dot x/\dot y=\pm py/\sqrt{1-p^2y^2}$, which integrates to
$x-x_0=\mp\frac1p\sqrt{1-p^2y^2}$, that is
\begin{equation}\label{eq:torus2:semicircle}
 (x-x_0)^2+y^2=\frac1{p^2}:
\end{equation}
a semicircle centred on the real axis, of radius $1/|p|$. An arc-length parametrization is
$x=x_0+r\tanh s$, $y=r/\cosh s$ (symbolic check: it solves the Euler--Lagrange equations
$\ddot x=2\dot x\dot y/y$, $\ddot y=(\dot y^2-\dot x^2)/y$ with unit speed). The arc of the
unit circle bounding $\mathcal F$ is a geodesic, and so are its vertical sides.

\begin{proposition}[Distance formula]\label{thm:torus2:distance}
The geodesic distance between $z,w\in\mathbb H$ is given by
\begin{equation}\label{eq:torus2:distance}
 \cosh d(z,w)=1+\frac{|z-w|^2}{2\operatorname{Im}z\,\operatorname{Im}w}.
\end{equation}
\end{proposition}
\begin{proof}
The right side is invariant under $SL(2,\R)$: from
$\gamma z-\gamma w=(z-w)/\big((cz+d)(cw+d)\big)$ and \eqref{eq:torus2:Im} the factors
$|cz+d|^2|cw+d|^2$ cancel. The left side is invariant because the metric is. By
transitivity we may take $z=\ii$; a rotation about $\ii$ then turns the geodesic through
$\ii$ and $w$ into the vertical one, so that $w=\ii\ee^{d}$. There the right side is
$1+(\ee^d-1)^2/(2\ee^d)=\frac12(\ee^d+\ee^{-d})$.
\end{proof}

\begin{physicsbox}[Physical picture: infinite distance and light towers]
Take a rectangular torus without $B$-field and let its area grow at fixed shape,
$\rho=\ii\rho_2\to\ii\infty$. By the case $p=0$ this is a geodesic, and the distance travelled
is $\Delta=\ln(\rho_2^{\rm final}/\rho_2^{\rm initial})$: the decompactification limit is
at \emph{infinite} distance\index{moduli space!distance}, although the fundamental domain has finite area. Along this
path the momentum states of \eqref{eq:torus2:pLpR} have
$M^2\propto|n_1-\tau n_2|^2/(\rho_2\tau_2)$, so a whole tower of states becomes light as
\[
 M\;\propto\;\rho_2^{-1/2}=\ee^{-\Delta/2}\,M_{\rm initial},
\]
exponentially in the distance (the coefficient $\frac12$ depends on the normalization
chosen in \eqref{eq:torus2:modmetric}). In the opposite direction $\rho_2\to0$ the winding
states do the same, as they must by $S'$. The cusp $\tau\to\ii\infty$ is the limit of a
long thin torus, with one light momentum tower and one light winding tower. That infinite
distances in moduli space come with exponentially light towers is the pattern elevated to
a conjecture about quantum gravity in \cref{ch:swampland}; on $T^2$ it is a computation.
For the circle, the combination $R+\ap/R$ of \cref{ch:dualscale} equals
$2\sqrt{\ap}\cosh u$ with $u=\ln(R/\sqrt{\ap})$ proportional to the distance from the
self-dual point: it is minimal there and grows exponentially in both directions.
\end{physicsbox}

\section{Mirror symmetry as a factorized duality}\label{sec:torus2:mirror}

For Calabi--Yau manifolds, mirror symmetry\index{mirror symmetry} is a duality that
``relates the complex-structure deformations of one CY background to the K\"ahler-class
deformations of its mirror partner'' \source{GPR §1, ll.~334--340}; see also
\source{Asp §3.4, ll.~1840--1842}. The torus $T^2$ is the Calabi--Yau manifold of complex
dimension one, and for it the statement is a theorem we have already met:
\begin{equation}\label{eq:torus2:mirror}
 D_2:\qquad(\tau,\rho)\;\longleftrightarrow\;(\rho,\tau)\qquad\TL
\end{equation}
``For a complex torus, mirror symmetry is identical to what is called `factorized duality'\,''
\source{GPR §1, ll.~344--345}, and ``an $R\leftrightarrow1/R$ transformation on one of the
circles in the $T^2$ is a mirror map in the sense that the notions of deformation of complex
structure and complexified K\"ahler form are interchanged'' \source{Asp §5.1,
ll.~2826--2829}. The rectangular torus of \cref{ex:torus2:rect} shows the mechanism in one
line. Three consequences deserve emphasis.

\paragraph{The mirror manifold is a different manifold.} The torus with moduli
$(\tau,\rho)$ and the torus with moduli $(\rho,\tau)$ are different Riemannian manifolds:
in \cref{fig:torus2:lattice}b one is long and thin, the other large and nearly square. A
point particle distinguishes them. The string does not, because it has winding states in
addition to momentum states, and the duality exchanges $n_2\leftrightarrow m_2$.

\paragraph{Non-geometric symmetries are mirrors of geometric ones.} By
\cref{thm:torus2:group}, $T'=D_2TD_2$ and $S'=D_2SD_2$. The periodicity of the $B$-field is
the mirror image of the freedom to change the lattice basis by
$\hat e_1\to\hat e_1+\hat e_2$, and the inversion of the area is the mirror image of the
exchange of the two sides of the cell. In general, conjugation by $D_2$ is an isomorphism
$SL(2,\Z)_\tau\to SL(2,\Z)_\rho$. This is how mirror symmetry is used in practice, also for
Calabi--Yau threefolds: a quantity that is hard to compute on the K\"ahler side, where the
string corrections live, is computed on the complex-structure side of the mirror, where
classical geometry suffices \source{GPR §1, ll.~338--340}.

\paragraph{The other circle.} Had we dualized the first circle instead, \eqref{eq:torus2:D}
gives $(\tau,\rho)\to(-1/\rho,-1/\tau)$: still an exchange of shape and size, composed with
the two inversions. Conjugating the shift $T$ by $D_1$ we find, following a point
$(\tau,\rho)$ through the three maps of $g=D_1g_TD_1$ in the action \eqref{eq:torus2:action}
(first $D_1$, then $T$ on the first entry, then $D_1$),
\begin{equation}\label{eq:torus2:D1TD1}
 (\tau,\rho)\to\Big(-\frac1\rho,-\frac1\tau\Big)\to\Big(1-\frac1\rho,-\frac1\tau\Big)
 \to\Big(\tau,\ \frac{-1}{1-1/\rho}\Big)=\Big(\tau,\frac{\rho}{1-\rho}\Big).
\end{equation}
This is the element
$\big(\begin{smallmatrix}1&0\\-1&1\end{smallmatrix}\big)=S'T'S'^{-1}$ of $SL(2,\Z)_\rho$: a
shift by one unit not of $\rho$ but of $-1/\rho$. As an integer matrix it is the
\emph{lower} triangular
$\big(\begin{smallmatrix}\mathbf 1&0\\J&\mathbf 1\end{smallmatrix}\big)=g_{-J}^{\mathsf T}$.
This precise identity is what has been formalized, as \lean{mirror\_conjugates\_tauShift};
see \cref{sec:torus2:lean} for its reading in both conventions for the action.

\begin{remark}[Type II strings]
For the superstring, T-duality on one circle exchanges the type IIA and type IIB theories
\source{Asp §5.1, Proposition~5, ll.~2821--2825}. Hence type IIA on a torus with moduli
$(\tau,\rho)$ is type IIB on the torus with moduli $(\rho,\tau)$, and the modular group that
is geometric in one description acts on the K\"ahler modulus in the other
\source{Asp §5.1, ll.~2826--2831}. This is the form in which the statement is used for
$\KT$ in \cref{ch:k3t2}.
\end{remark}

\begin{remark}[A third half-plane]
The one-loop partition function depends on $\tau$, $\rho$ and the modulus $\sigma$ of the
worldsheet torus, and it is symmetric under the exchange of the target-space complex
structure with the worldsheet one, $Z(\tau,\rho;\sigma)=Z(\sigma,\rho;\tau)$; together with
$\tau\leftrightarrow\rho$ this is a triality \source{GPR §2.4.2, eq.~(2.4.62),
ll.~2005--2016} \TL. We do not use it in this book.
\end{remark}

\section{Worked examples}\label{sec:torus2:examples}

\begin{example}[A square torus in disguise]\label{ex:torus2:disguise}
A background is given by
\[
 G=\begin{pmatrix}13&8\\8&5\end{pmatrix},\qquad B_{12}=\frac73 .
\]
\emph{Moduli.} $\det G=65-64=1$, so by \eqref{eq:torus2:tau} and \eqref{eq:torus2:rho}
$\tau=\frac85+\frac15\ii$ and $\rho=\frac73+\ii$. The cell has the area of the self-dual
square torus, but it looks very elongated: $|\tau|^2=\frac{65}{25}$, and the angle between
the sides has $\cos\theta=\tau_1/|\tau|=8/\sqrt{65}\approx0.992$, $\theta\approx7^\circ$.

\emph{Reduction of $\tau$.} Step 1: $T^{-2}$ gives $\tau-2=-\frac25+\frac15\ii$. Step 2:
$|\tau-2|^2=\frac4{25}+\frac1{25}=\frac15<1$, so apply $S$:
$-1/z=-\bar z/|z|^2=5\,(\frac25+\frac15\ii)=2+\ii$. Step 1 again: $T^{-2}$ gives $\ii$. Hence
$\tau\sim\ii$: the lattice is the square lattice in a skew basis. Indeed
$G=AA^{\mathsf T}$ with $A=\big(\begin{smallmatrix}2&3\\1&2\end{smallmatrix}\big)$,
$\det A=1$, and \eqref{eq:torus2:mobius} gives $(2\ii+3)/(\ii+2)=\frac15(3+2\ii)(2-\ii)=
\frac15(8+\ii)$.

\emph{Reduction of $\rho$.} $T'^{-2}$ gives $\rho-2=\frac13+\ii$, which lies in $\mathcal F$.
The background is equivalent to the square torus of unit area with $B_{12}=\frac13$. It is
\emph{not} equivalent to the $SU(2)^2$ point $(\ii,\ii)$. The orbit of $(\ii,\ii)$ under
the whole group consists of pairs $(\gamma\ii,\gamma'\ii)$, because $D_2$ and $R$ map
$(\ii,\ii)$ to itself; and $\frac13+\ii$ is an interior point of $\mathcal F$, hence not
equivalent to the point $\ii$ of $\mathcal F$.

\emph{Distance to the enhanced point.} By \eqref{eq:torus2:distance},
$\cosh d=1+\frac{(1/3)^2}{2}=\frac{19}{18}$, $d\approx0.332$.
\end{example}

\begin{example}[The two enhanced symmetry points]\label{ex:torus2:enhanced}
GPR give the $SU(3)$ point as $E=\big(\begin{smallmatrix}1&1\\0&1\end{smallmatrix}\big)$
\source{GPR §2.4.2, eq.~(2.4.64), ll.~2026--2036}. Its symmetric and antisymmetric parts are
$G=\big(\begin{smallmatrix}1&1/2\\1/2&1\end{smallmatrix}\big)$ and $B_{12}=\frac12$, so
$\det G=\frac34$ and
\[
 \tau=\frac{1/2}{1}+\ii\frac{\sqrt{3}/2}{1}=\ee^{\ii\pi/3},\qquad
 \rho=\frac12+\ii\frac{\sqrt3}2=\ee^{\ii\pi/3},
\]
in agreement with \eqref{eq:torus2:special}. The point is fixed by $TS$ acting on $\tau$, by
$T'S'$ acting on $\rho$, and by the mirror map $D_2$, because $\tau=\rho$. As a check on
\eqref{eq:torus2:pLpR}, take $(n_1,n_2;m_1,m_2)=(1,0;1,0)$: $a=1$, $c=\tau$, and with
$\rho=\tau=\omega:=\ee^{\ii\pi/3}$, $2\rho_2\tau_2=\frac32$,
\[
 p_L^2=\tfrac23\,|1-\omega^2|^2=\tfrac23\cdot3=2,\qquad
 p_R^2=\tfrac23\,|1-\bar\omega\omega|^2=0 ,
\]
using $\omega^2=\ee^{2\pi\ii/3}$ and $|1-\ee^{2\pi\ii/3}|^2=3$. With $N=1,\tilde N=0$,
\eqref{eq:torus2:mass} gives $M^2=0$: one of the six charged massless vectors of $SU(3)_L$
found in \cref{ch:narain}. A finite search over $|n_i|,|m_i|\le2$ (symbolic check) finds
six vectors with $(p_L^2,p_R^2)=(2,0)$ at this point, four at $(\ii,\ii)$ and none at
$(\tau,\rho)=(\ii,2\ii)$.

How far apart are the two enhanced points? In each half-plane,
\eqref{eq:torus2:distance} with $z=\ii$, $w=\omega$ gives
$|z-w|^2=\frac14+(1-\frac{\sqrt3}2)^2=2-\sqrt3$ and
\[
 \cosh d=1+\frac{2-\sqrt3}{\sqrt3}=\frac2{\sqrt3},\qquad
 d=\ln\frac{2+1}{\sqrt3}=\tfrac12\ln3\approx0.549 ,
\]
This is also the distance in the quotient $\SLZ\backslash\mathbb H$: the images of
$\omega$ are the points $\omega+k$, $k\in\Z$, of which $\omega$ and $\omega-1$ are nearest
to $\ii$, and points with $\operatorname{Im}\le\frac{\sqrt3}{2}/3$ (because
$|c\omega+d|^2=c^2+cd+d^2$ is $1$ or at least $3$), which are at distance at least
$\ln(2\sqrt3)\approx1.24$ from $\ii$. Since both moduli move, the distance between the two
enhanced points in the metric \eqref{eq:torus2:modmetric} is $\sqrt2\,d\approx0.777$.
\end{example}

\section{What the machine has checked}\label{sec:torus2:lean}

Two files in the directory \texttt{DualScaleStream2/TDuality} treat $d=2$, namely
\texttt{Mirror.lean} and \texttt{SL2Product.lean}. They work entirely with $4\times4$
\emph{integer} matrices, built from three definitions valid for every $d$ (\cref{ch:odd}):
\lean{thetaShift}, which maps $\Theta$ to
$\big(\begin{smallmatrix}1&\Theta\\0&1\end{smallmatrix}\big)$; \lean{basisChange}, which
maps $A,B$ to $\diag(A,B)$; and \lean{factorized}, with \lean{factorized k} equal to
$D_{k+1}$ (Lean counts from $0$, so \lean{factorized 0} is $D_1$). Neither file mentions $\tau$ or
$\rho$: the complex moduli, the half-plane and \cref{thm:torus2:spectrum} are not
formalized. What is proved are matrix identities in the charge lattice, whose translation
into statements about moduli is the (unformalized, symbolic-check) dictionary of
\cref{sec:torus2:group}.

\begin{leanbox}[In Lean: the mirror conjugate of the $\tau$-shift (\texttt{Mirror.lean})]
\begin{leancode}
def tauShift : Matrix (Fin 2) (Fin 2) ℤ := !![1, 1; 0, 1]
def tauShiftDual : Matrix (Fin 2) (Fin 2) ℤ := !![1, 0; -1, 1]
theorem tauShift_dual_spec : tauShiftᵀ * tauShiftDual = 1 := by ...
theorem tauShift_isODD : IsODD (basisChange tauShift tauShiftDual) := ...
def mirrorTheta : Matrix (Fin 2) (Fin 2) ℤ := !![0, -1; 1, 0]
theorem mirrorTheta_antisymm : mirrorThetaᵀ = -mirrorTheta := by ...
theorem mirror_conjugates_tauShift :
    factorized (0 : Fin 2) * basisChange tauShift tauShiftDual * factorized 0 =
      (thetaShift mirrorTheta)ᵀ := by ...
\end{leancode}
\textbf{Reading.} With $T=\big(\begin{smallmatrix}1&1\\0&1\end{smallmatrix}\big)$ and
$\Theta=-J$: the pair $(T,(T^{\mathsf T})^{-1})$ satisfies the compatibility condition, so
$g_T\in\Odd{2}$; $\Theta$ is antisymmetric; and
$D_1\,g_T\,D_1=g_{-J}^{\mathsf T}$, the lower triangular matrix of
\eqref{eq:torus2:D1TD1}. This is one identity between two fixed $4\times4$ integer matrices.
\textbf{Proof idea.} \lean{ext} reduces the matrix equation to its 16 entries,
\lean{fin\_cases} enumerates them, and each is closed by \lean{rfl}: the kernel evaluates
both sides. \textbf{What it does not say.} Nothing about $\tau$, $\rho$ or the generalized
metric, and nothing about $D_2$; the companion $D_2\,g_T\,D_2=g_J$, which is GPR's
$T'=D_2TD_2$, is a symbolic check in this book (\cref{ex:torus2:lean2}).
\end{leanbox}

\begin{pitfallbox}[Common pitfall: one identity, two readings]
How \lean{mirror\_conjugates\_tauShift} translates to moduli depends on the convention for
the action. With $\HH\to g\HH g^{\mathsf T}$ (this chapter, GPR, \lean{genMetric\_bshift}),
$g_T$ is $\tau\to\tau+1$ and the right side is $\rho\to\rho/(1-\rho)$, as in
\eqref{eq:torus2:D1TD1}. With $\HH\to g^{\mathsf T}\HH g$ (\cref{ch:narain},
\lean{massForm\_covariant}), the right side is the plain $B$-shift $\rho\to\rho-1$, and
$g_T$ acts as $G\to T^{\mathsf T}GT$, which is $\tau\to\tau/(\tau+1)$. The docstring of the
Lean file describes the left side in the first reading and the right side in the second; it
also states explicitly that no sign or orientation of the $\rho$-shift is claimed. In
either convention the content is the same: $D_1$ conjugates a parabolic generator of
$SL(2,\Z)_\tau$ into a parabolic generator of $SL(2,\Z)_\rho$. Both readings have been
confirmed symbolically.
\end{pitfallbox}

\begin{leanbox}[In Lean: $SL(2,\Z)_\tau$ commutes with the $\rho$-translations
(\texttt{SL2Product.lean})]
\begin{leancode}
def jMat : Matrix (Fin 2) (Fin 2) ℤ := !![0, 1; -1, 0]
theorem mul_jMat_mul_transpose (A : Matrix (Fin 2) (Fin 2) ℤ) :
    A * jMat * Aᵀ = A.det • jMat := by ...
theorem basisChange_mul {d : ℕ} (A B A' B' : Matrix (Fin d) (Fin d) ℤ) :
    basisChange A B * basisChange A' B' = basisChange (A * A') (B * B') := by ...
theorem thetaShift_mul {d : ℕ} (Θ Θ' : Matrix (Fin d) (Fin d) ℤ) :
    thetaShift Θ * thetaShift Θ' = thetaShift (Θ + Θ') := by ...
theorem basisChange_comm_thetaShift (A B : Matrix (Fin 2) (Fin 2) ℤ)
    (hAB : Aᵀ * B = 1) (hdet : A.det = 1) (t : ℤ) :
    basisChange A B * thetaShift (t • jMat) =
      thetaShift (t • jMat) * basisChange A B := by ...
\end{leancode}
\textbf{Reading.} The first theorem is \eqref{eq:torus2:AJAT}, for all integer $2\times2$
matrices. The next two say that basis changes multiply blockwise and that $\Theta$-shifts
add, $g_\Theta g_{\Theta'}=g_{\Theta+\Theta'}$, for every $d$. The last is
\eqref{eq:torus2:commute}: \emph{if} $A^{\mathsf T}B=1$ and $\det A=1$, \emph{then}
$g_A$ commutes with $T'^{\,t}=g_{tJ}$ for every integer $t$. \textbf{Proof idea.} From
$\det A=1$ and the first theorem, $AJA^{\mathsf T}=J$; multiplying $AJ$ by
$\mathbf 1=A^{\mathsf T}B$ on the right and regrouping gives $AJ=JB$; then both products
are expanded blockwise and compared entry by entry. \textbf{What it does not say.} Only
one direction is formalized. The converse --- for $t\ne0$ the two matrices commute
\emph{only if} $\det A=1$ --- follows on paper from
$AJ-JA^{-\mathsf T}=(1-1/\det A)\,AJ$ and has been confirmed symbolically, but is not in
Lean. The inversion $S'$ and hence the full $SL(2,\Z)_\rho$ do not appear, nor does
\cref{thm:torus2:group}.
\end{leanbox}

The half-plane geometry of \cref{sec:torus2:halfplane} is touched by two files of an older library of this project
(\texttt{StringTheoryFormalization}), written for other purposes.

\begin{leanbox}[In Lean: composition of M\"obius maps (\texttt{Frontier/SL2CSymmetry.lean})]
\begin{leancode}
structure MobiusTransform where
  (a b c d : ℂ)
  det_one : a * d - b * c = 1
noncomputable def MobiusTransform.act (M : MobiusTransform) (z : ℂ) : ℂ :=
  (M.a * z + M.b) / (M.c * z + M.d)
theorem mobius_compose (M N : MobiusTransform) (z : ℂ)
    (h : N.c * z + N.d ≠ 0)
    (h' : M.c * (N.act z) + M.d ≠ 0) :
    (MobiusTransform.mk
      (M.a * N.a + M.b * N.c)
      (M.a * N.b + M.b * N.d)
      (M.c * N.a + M.d * N.c)
      (M.c * N.b + M.d * N.d)
      (by
        have key : (M.a * N.a + M.b * N.c) * (M.c * N.b + M.d * N.d)
            - (M.a * N.b + M.b * N.d) * (M.c * N.a + M.d * N.c)
            = (M.a * M.d - M.b * M.c) * (N.a * N.d - N.b * N.c) := by ring
        rw [key, M.det_one, N.det_one, one_mul])).act z =
    M.act (N.act z) := by ...
theorem mobius_id_act (z : ℂ) : mobiusId.act z = z := by ...
\end{leancode}
\textbf{Reading.} For complex matrices of determinant one, the matrix product (whose
determinant is shown to be one inside the statement) acts as the composition of the two
maps, at every $z$ where the two denominators do not vanish; the identity matrix acts
trivially. This is \eqref{eq:torus2:compose} over $\C$; it contains the real and the integer
case. \textbf{Proof idea.} Unfold, clear denominators with \lean{field\_simp}, finish with
\lean{ring}. \textbf{What it does not say.} That $\mathbb H$ is preserved
\eqref{eq:torus2:Im}, anything about $\SLZ$ or a fundamental domain (for these see Mathlib's
\lean{ModularGroup}), or about the metric. The file was written for the conformal group of
the sphere (\cref{ch:cft}); its theorems named \lean{ward\_identity\_translation} and
\lean{ward\_identity\_dilatation} have the conclusion \lean{True} and prove nothing.
\end{leanbox}

\begin{leanbox}[In Lean: what exists about geodesics (\texttt{Frontier/ModuliGeodesics.lean})]
\begin{leancode}
theorem poincare_geodesic_kinetic_energy_nonneg (y y' : ℝ) (hy : 0 < y) :
    let energy := (y')^2 / y^2
    0 ≤ energy := by ...
\end{leancode}
\textbf{Reading.} For real $y>0$ and any real $y'$, the number $y'^2/y^2$ is non-negative:
the $y$-part of the conserved energy in \eqref{eq:torus2:conserved} is a square. It is a
statement about two real numbers, proved by \lean{positivity}; no metric, geodesic equation
or conservation law is formalized. The other declarations of this file must be read with
the same care, as its own documentation says: \lean{poincare\_geodesic\_atlas\_curvature}
checks by \lean{rfl} that a structure field whose default value is $-1$ has the value $-1$
--- it is not the computation \eqref{eq:torus2:K}; \lean{geodesic\_equation\_kummer\_locus}
has a conclusion of the form ``$\ldots\vee\lean{True}$''; and
\lean{wp\_geodesic\_completeness} is proved by extending a path by itself. None of them is
evidence for a statement of this chapter.
\end{leanbox}

\begin{table}[ht]
\centering\small
\begin{tabular}{@{}p{0.50\textwidth}cp{0.37\textwidth}@{}}\toprule
Claim & Tier & Evidence\\\midrule
$AJA^{\mathsf T}=(\det A)J$ for integer $2\times2$ $A$ & \TA & \lean{mul\_jMat\_mul\_transpose}\\
$\det A=1\Rightarrow g_A$ commutes with $g_{tJ}$ & \TA & \lean{basisChange\_comm\_thetaShift}\\
Converse of the previous line & --- & on the page; symbolic check\\
$g_Ag_{A'}=g_{AA'}$, $g_\Theta g_{\Theta'}=g_{\Theta+\Theta'}$ (all $d$) & \TA &
\lean{basisChange\_mul}, \lean{thetaShift\_mul}\\
$g_T\in\Odd2$; $D_1g_TD_1=g_{-J}^{\mathsf T}$; $D_1D_2=\eta$ & \TA & \lean{tauShift\_isODD},
\lean{mirror\_conjugates\_tauShift}, \lean{factorized\_two}\\
$g_\Theta\HH(G,B)g_\Theta^{\mathsf T}=\HH(G,B+\Theta)$ (real matrices) & \TA &
\lean{genMetric\_bshift}\\
M\"obius maps compose like matrices (over $\C$) & \TA & \lean{mobius\_compose}\\
$p_{L,R}^2$ in terms of $\tau,\rho$, \eqref{eq:torus2:pLpR} & \TL & GPR (2.4.58); derivation
and symbolic check\\
Action of $S,T,S',T',D_1,D_2,R,W$ on $(\tau,\rho)$ & \TL & GPR (2.4.57)--(2.4.61); symbolic
check ($D_1$: this book)\\
$\Odd2\cong[SL(2,\Z)\times SL(2,\Z)]\rtimes(\Z_2\times\Z_2)$ & \TL & GPR §2.4.2\\
Mirror symmetry of $T^2$ $=$ factorized duality & \TL & GPR §1; Asp §5.1\\
Metric \eqref{eq:torus2:modmetric}, curvature, geodesics, distance formula & --- & derived on
the page; symbolic checks\\
\bottomrule
\end{tabular}
\caption{Epistemic status of the claims of this chapter. The identification of the integer
matrices with dualities of a physical string background is never Tier A. All Tier A entries
depend only on the axioms \lean{propext}, \lean{Classical.choice} and \lean{Quot.sound}.}
\label{tab:torus2:tiers}
\end{table}

\paragraph{What a faithful formalization would need.} (i) A definition of $\tau$ and $\rho$
as points of Mathlib's \lean{UpperHalfPlane} from a positive definite real $G$ and an
antisymmetric $B$, with the bijection \eqref{eq:torus2:E}. (ii) The theorem that
$g_A\HH g_A^{\mathsf T}$ has moduli $(A\cdot\tau,\rho)$ for $A\in\SLZ$, where Mathlib's
action of $\SLZ$ on the half-plane can be reused, and the analogous statements for
$g_{tJ}$, $D_1$, $D_2$. (iii) \Cref{thm:torus2:spectrum} as an identity of real polynomials
divided by $\rho_2\tau_2$ --- within reach of \lean{field\_simp} and \lean{ring}. (iv) The
structure theorem \eqref{eq:torus2:structure}, which needs $\Odd2$ as a group and a
generation argument; this is the hard part.

\begin{summarybox}
\begin{itemize}[nosep,leftmargin=*]
\item The moduli of $T^2$ are $\tau=(G_{12}+\ii\sqrt{\det G})/G_{22}$, the shape, and
$\rho=B_{12}+\ii\sqrt{\det G}$, the area in string units with the $B$-field; the moduli space
is locally $\mathbb H\times\mathbb H$.
\item $p_L^2$ and $p_R^2$ are $|(n_1-\tau n_2)-\rho(m_2+\tau m_1)|^2/(2\rho_2\tau_2)$ with
$\rho$, respectively $\bar\rho$; dualities are the changes of $(\tau,\rho)$ that can be undone
by relabelling the integers.
\item $SL(2,\Z)_\tau$ (basis changes) is geometric; $SL(2,\Z)_\rho$ is generated by the
integer $B$-shift $T'$ and the area inversion $S'$. The two commute; together with the
exchange $D_2$ and the reflection $R$ they generate $\Odd2$.
\item Mirror symmetry of $T^2$ is T-duality on one circle, $\tau\leftrightarrow\rho$; it maps
the geometric modular group to the stringy one, $T'=D_2TD_2$, $S'=D_2SD_2$.
\item Each half-plane carries the Poincar\'e metric, of curvature $-1$; geodesics are
vertical lines and semicircles; $\cosh d=1+|z-w|^2/(2\operatorname{Im}z\operatorname{Im}w)$.
Decompactification is at infinite distance, with a tower of mass $\propto\ee^{-\Delta/2}$.
\item Kernel-checked: $AJA^{\mathsf T}=(\det A)J$, the commutation of $g_A$ ($\det A=1$) with
$g_{tJ}$ (one direction), $D_1g_TD_1=g_{-J}^{\mathsf T}$, $D_1D_2=\eta$, composition of M\"obius
maps. Not formalized: $\tau$, $\rho$, their transformation laws, the group structure, the
hyperbolic geometry.
\end{itemize}
\end{summarybox}

\section*{Exercises}
\addcontentsline{toc}{section}{Exercises}

\begin{exercise}[$\star$]
For $G=\big(\begin{smallmatrix}5&3\\3&2\end{smallmatrix}\big)$ and $B_{12}=\frac52$ compute
$\tau$ and $\rho$, and reduce both to $\mathcal F$. Why is the reduced $\rho$ not unique?
(Answer: $\tau\sim\ii$; $\rho\sim\pm\frac12+\ii$, two boundary points identified by $T'$.)
\end{exercise}

\begin{exercise}[$\star$]
Verify \eqref{eq:torus2:Im} and \eqref{eq:torus2:compose} by direct computation, and show that
$\gamma z-\gamma w=(z-w)/\big((cz+d)(cw+d)\big)$.
\end{exercise}

\begin{exercise}[$\star\star$]
Using \eqref{eq:torus2:pLpR}, show that $W:(\tau,\rho)\to(\tau,-\bar\rho)$ exchanges
$p_L^2$ and $p_R^2$ when combined with $m_i\to-m_i$. Conclude from $p_L^2-p_R^2=2n\cdot m$
that $W$ cannot be represented by an element of $\Odd2$ acting on $Z$.
\end{exercise}

\begin{exercise}[$\star\star$]
Derive the action of $D_1$ in \eqref{eq:torus2:D} from \eqref{eq:torus2:pLpR}: show that
the background $(-1/\rho,-1/\tau)$ with $n_1$ and $m_1$ exchanged has the same $p_L^2$ and
$p_R^2$ as $(\tau,\rho)$. (Hint: multiply the numerator by $|\tau\rho|^{-1}$.)
\end{exercise}

\begin{exercise}[$\star\star$]
Show that the backgrounds $E(g,b)=\big(\begin{smallmatrix}g&b\\-b&g\end{smallmatrix}\big)$
with $g^2+b^2=1$ have $\tau=\ii$ and $|\rho|=1$, and that $E^{-1}=E(g,-b)$. Which element of
the full symmetry group fixes every point of this circle?
\source{GPR §2.4.2, eqs.~(2.4.66)--(2.4.67), ll.~2053--2075}
\end{exercise}

\begin{exercise}[$\star\star$]
Find the geodesic through $z=\ii$ and $w=1+2\ii$ (centre and radius of the semicircle) and
compute $d(z,w)$ from \eqref{eq:torus2:distance}. (Answer: centre $x_0=2$, radius
$\sqrt5$; $\cosh d=\frac32$.)
\end{exercise}

\begin{exercise}[$\star\star\star$]
The domain $\mathcal F$ is a hyperbolic triangle with angles $\frac\pi3,\frac\pi3,0$. The
Gauss--Bonnet theorem gives the area of a geodesic triangle of curvature $-1$ as $\pi$ minus
the sum of its angles. Compare with \eqref{eq:torus2:area}, and compute the area of the
fundamental domain of the subgroup generated by $T^2$ and $S$.
\end{exercise}

\begin{exercise}[$\star$, Lean]
State and prove in Lean that $J^2=-\mathbf 1$ (for \lean{jMat}) and that the square of
\lean{tauShift} is \lean{!![1, 2; 0, 1]}. The proof pattern of
\lean{mirrorTheta\_antisymm} suffices: \lean{ext}, then \lean{fin\_cases}, then \lean{rfl}
or \lean{decide}.
\end{exercise}

\begin{exercise}[$\star\star$, Lean]\label{ex:torus2:lean2}
State and prove the $D_2$ companion of \lean{mirror\_conjugates\_tauShift}:
\begin{leancode}
example : factorized (1 : Fin 2) * basisChange tauShift tauShiftDual * factorized 1
    = thetaShift jMat := by sorry
\end{leancode}
This is GPR's $T'=D_2TD_2$. Then use \lean{factorized\_mul\_self} to deduce
$g_T=D_2\,g_J\,D_2$.
\end{exercise}

\begin{exercise}[$\star\star\star$, Lean]
Formulate the converse of \lean{basisChange\_comm\_thetaShift}: if $A^{\mathsf T}B=1$,
$t\neq0$ and the two matrices commute, then $\det A=1$. (Hint: compare the upper right
blocks to get $t\,AJ=t\,JB$, multiply by $A^{\mathsf T}$ on the right, and use
\lean{mul\_jMat\_mul\_transpose}; you will need that $J\ne0$ and that $\Z$ has no zero
divisors.)
\end{exercise}

\section*{Further reading}
\addcontentsline{toc}{section}{Further reading}
\begin{itemize}[nosep,leftmargin=*]
\item The $d=2$ example, with the generators $S,T,S',T',D_2,R,W$, the fixed points and the
triality with the worldsheet modulus: \source{GPR §2.4.2, ll.~1874--2075}. The generators of
$\Odd d$ for general $d$: \source{GPR §2.4, ll.~1340--1423}.
\item Mirror symmetry as factorized duality, and its place among target-space dualities:
\source{GPR §1, ll.~334--348}.
\item The moduli space of conformal field theories on $T^2$ as a product of two copies of
$SL(2)/U(1)$, and T-duality on one circle as the mirror map between type IIA and type IIB:
\source{Asp §5.1, ll.~2750--2762 and 2819--2831}. Mirror symmetry for K3, where the
distinction between complex structure and K\"ahler form is more subtle:
\source{Asp §3.4, ll.~1836--1888}, taken up in \cref{ch:stringsk3}.
\item The generalized-metric form of the kinetic term: \source{HHZ §1, ll.~334--346};
\cref{ch:genmetric,ch:dftaction}.
\item Lean sources: the files \texttt{Mirror.lean}, \texttt{SL2Product.lean},
\texttt{ODD.lean} and \texttt{Factorized.lean}, all in the directory
\texttt{DualScaleStream2/TDuality}. For the modular group and its fundamental domain see the Mathlib
file \texttt{NumberTheory/Modular.lean}.
\end{itemize}
